\section{Optional rational specialization: \texorpdfstring{$X^5-X-1$}{X5-X-1}}
\label{sec:explicit-quintic}

The generic proof is the cleanest route for the iPRA architecture because its
Galois action is literal variable permutation.  For readers who want a single
polynomial over $\Q$, this section gives a compact finite certificate for
\[
 f(X)=X^5-X-1.
 \tag{E1}\label{eq:selmer5}
\]
The only additional theorem needed is the finite Dedekind--Frobenius lemma:
for a prime not dividing the discriminant, the degrees of the distinct
irreducible factors modulo that prime occur as the cycle lengths of an element
of the rational Galois group.  That lemma can itself be arithmetized, but doing
so adds more local algebra than the generic proof requires.

\subsection{Irreducibility modulo three}

Modulo $3$, the values at $0,1,2$ are all $2$, so $f$ has no linear factor.
The monic irreducible quadratics over $\F_3$ are
\[
 X^2+1,
 \qquad X^2+X+2,
 \qquad X^2+2X+2.
\]
Polynomial division gives the respective remainders
\[
 -1,
 \qquad X-1,
 \qquad X-1.
\]
Thus none divides $f$.  A reducible polynomial of degree five has a factor of
degree one or two: among the degrees of its nonconstant factors, the smallest
is at most $2$.  Hence $f$ is irreducible over $\F_3$, and Gauss's lemma gives
irreducibility over $\Q$.

The same modular factorization is a single irreducible factor of degree five,
so the Dedekind--Frobenius lemma supplies a $5$-cycle in the Galois group.

\subsection{A transposition from reduction modulo two}

Modulo $2$ one has the exact factorization
\[
 X^5+X+1
 =(X^2+X+1)(X^3+X^2+1).
 \tag{E2}\label{eq:mod2-factor}
\]
The quadratic has no root in $\F_2$, and the cubic also has no root in
$\F_2$; hence both factors are irreducible.  A resultant computation gives
\[
 \operatorname{disc}(X^5-X-1)=2869=19\cdot151.
 \tag{E3}\label{eq:selmer-disc}
\]
In particular, neither $2$ nor $3$ divides the discriminant.  Reduction
modulo $2$ is therefore square-free, and the factor degrees $2$ and $3$ give
an element of cycle type $(2)(3)$ in the Galois group.  The cube of this
element is a transposition.

\subsection{The Galois group is \texorpdfstring{$S_5$}{S5}}

Let $c$ be the $5$-cycle obtained modulo $3$ and let $t$ be the transposition
obtained from the cube of the $(2)(3)$ element.  The conjugates
\[
 c^k t c^{-k},\qquad k=0,1,2,3,4,
\]
are transpositions whose edges form a connected graph on the five letters.
Indeed, if $t$ exchanges vertices at cyclic distance $d$, then the conjugates
join $i$ to $i+d$ modulo $5$; because $5$ is prime and $d\ne0$, this graph is
a $5$-cycle.  Transpositions along the edges of any connected graph generate
the full symmetric group.  Hence
\[
 \Gal(f/\Q)=S_5.
\]
By Theorem~\ref{thm:S5-not-solvable}, this group is not solvable.

\begin{corollary}[Explicit rational Abel--Ruffini witness]
No root of $X^5-X-1$ is expressible by radicals over $\Q$.
\end{corollary}

\begin{proof}
Suppose one root $\alpha$ lay in a radical tower.  Take a finite solvable
normal closure $N/\Q$ of that tower as in
Theorem~\ref{thm:radical-group-solvable}.  Since $f$ is irreducible, every
conjugate root is the image of $\alpha$ under a $\Q$-embedding; normality of
$N$ places all those images in $N$.  Thus the splitting field of $f$ is
contained in $N$, and its Galois group is a quotient of the solvable group
$\Gal(N/\Q)$.  This contradicts $\Gal(f/\Q)=S_5$.
\end{proof}

This rational example strengthens the coefficient-field interpretation of
Section~\ref{sec:abel-proof}.  It is not needed for the core generic proof, but
it is useful as an independently checkable specialization and aligns with the
rootwise rational results already developed in ProveIt.
